% ============================================================
%  PREÀMBUL — Dossier: Pneumàtica (4t ESO)
%  Mateix estil que «Circuits amb resistències».
%  Els símbols pneumàtics són a  pneumatica-tikz.sty
% ============================================================
\usepackage[utf8]{inputenc}
\usepackage[T1]{fontenc}
\usepackage{textcomp}
% El punt volat (·) es dibuixa amb el símbol matemàtic \cdot: així es veu igual
% a qualsevol tipus de lletra i amb qualsevol motor (pdfLaTeX, XeTeX/Tectonic, LuaTeX).
\usepackage{newunicodechar}
\newunicodechar{·}{\ensuremath{\cdot}}
\newunicodechar{«}{\guillemotleft}
\newunicodechar{»}{\guillemotright}
\newcommand{\sep}{\ \ensuremath{\cdot}\ }
\usepackage[catalan]{babel}
\usepackage[sc]{mathpazo}
\usepackage[a4paper,left=2cm,right=2cm,top=2.3cm,bottom=2.2cm]{geometry}
\usepackage{microtype}
\usepackage{xcolor}
\usepackage{amsmath,amssymb}
\usepackage{siunitx}
\sisetup{output-decimal-marker={,},group-minimum-digits=5,group-separator={\,},
         detect-all=false,per-mode=symbol,range-phrase={ a },range-units=single}
\DeclareSIUnit{\atm}{atm}
\usepackage{iftex}
\usepackage{tikz}
\usetikzlibrary{arrows.meta,calc,positioning,shapes.geometric,patterns,decorations.pathreplacing}
\usepackage{pneumatica-tikz}
\usepackage[most]{tcolorbox}
\usepackage{booktabs,array,enumitem,fancyhdr,titlesec,needspace,pifont,graphicx}
\usepackage[hidelinks]{hyperref}

\linespread{1.08}
\setlength{\parindent}{0pt}
\setlength{\parskip}{0.5em}
\setlist{nosep,leftmargin=1.4em,itemsep=0.3em}

% ---------- Colors ----------
\definecolor{blau}{HTML}{1F5F99}
\definecolor{verd}{HTML}{2E7D32}
\definecolor{taronja}{HTML}{E07B00}
\definecolor{vermell}{HTML}{B3261E}
\definecolor{lila}{HTML}{6A3FA0}
\definecolor{aire}{HTML}{0B7C8C}% color de l'aire comprimit (tubs, fletxes)

% ---------- Títols ----------
\titleformat{\section}{\Large\bfseries\color{blau}}{Bloc \thesection}{0.7em}{}[\vspace{-0.4em}\color{blau}\titlerule]
\titleformat{\subsection}{\large\bfseries\color{blau}}{}{0em}{}
\titlespacing*{\section}{0pt}{0pt}{1em}
\titlespacing*{\subsection}{0pt}{1.2em}{0.4em}

% ---------- Capçalera i peu ----------
\pagestyle{fancy}
\fancyhf{}
\renewcommand{\headrulewidth}{0.3pt}
\fancyhead[L]{\small\textsc{Pneumàtica}\sep 4t ESO}
\fancyhead[R]{\small Dossier de pneumàtica\ifsolucions\sep\textbf{SOLUCIONARI}\fi}
\fancyfoot[C]{\small\thepage}

% ---------- Quadres de colors ----------
\tcbset{boxrule=0.6pt, arc=3pt, left=7pt, right=7pt, top=5pt, bottom=5pt,
        fonttitle=\bfseries, breakable, before skip=0.8em, after skip=0.8em}
\newtcolorbox{idea}{colback=blau!6, colframe=blau, title={La idea clau}}
\newtcolorbox{metode}[1]{colback=lila!6, colframe=lila, title={Mètode: #1}}
\newtcolorbox{exemple}[1]{colback=verd!7, colframe=verd, title={Exemple resolt: #1}}
\newtcolorbox{guiada}[1]{colback=taronja!8, colframe=taronja, title={Practiquem junts: #1}}
\newtcolorbox{atencio}{colback=vermell!6, colframe=vermell, title={Compte amb l'error típic}}
\newtcolorbox{comprova}{colback=yellow!12, colframe=yellow!55!black, title={Comprova-ho}}
\newtcolorbox{sabiesque}{colback=aire!7, colframe=aire, title={Ho sabies?}}

% ---------- Nivells d'exercici ----------
\newcommand{\estrella}{\ding{72}}
\newcommand{\nivell}[1]{%
  \ifcase#1\or{\color{verd}\estrella}\or{\color{taronja}\estrella\estrella}\or{\color{vermell}\estrella\estrella\estrella}\fi}

\newcounter{exercici}[section]
\renewcommand{\theexercici}{\thesection.\arabic{exercici}}
\newenvironment{exercici}[2][8]{%
  \par\bigskip\needspace{#1\baselineskip}\noindent
  \refstepcounter{exercici}\textbf{\theexercici}\enspace\nivell{#2}\enspace\ignorespaces}%
  {\par}

% ---------- Solucions (espai en blanc o quadre de solució) ----------
%  \resposta{espai en blanc}{solució}
\newcommand{\resposta}[2]{%
  \ifsolucions
    \par\smallskip
    \begin{tcolorbox}[colback=verd!6, colframe=verd!60!black, boxrule=0.5pt, arc=2pt,
                      left=5pt, right=5pt, top=3pt, bottom=3pt, before skip=0.3em, after skip=0.6em]
      {\small\textbf{Solució.}\ #2}
    \end{tcolorbox}
  \else
    \vspace*{#1}
  \fi}

% ---------- Passos del mètode (per a les pràctiques guiades) ----------
\newcommand{\pas}[3]{\par\noindent\textbf{\textcolor{taronja}{\ding{\numexpr171+#1\relax}~#2}}\quad #3\par\smallskip}

% ---------- Dibuixos reutilitzables ----------
% Símbol dins del text o d'una taula: \simb{ <ordres de pneumatica-tikz> }
\newcommand{\simb}[2][0pt]{\raisebox{#1}{\begin{tikzpicture}[baseline=(current bounding box.center)]#2\end{tikzpicture}}}

% Diagrama espai-fase (1 = tija fora, 0 = tija dins; cal "fases+1" valors per cilindre)
%   \espaifaseu{fases}{A}{0,1,1,0}              (un cilindre)
%   \espaifase{fases}{A}{0,1,1,0}{B}{0,0,1,0}   (dos cilindres)
\newcommand{\efframe}[2]{%
  \foreach \k in {0,...,#1}{\draw[gray!35] (\k,-.25)--(\k,#2); \node[above,font=\scriptsize] at (\k,#2) {\k};}
  \node[above,font=\scriptsize\itshape] at ({#1/2},{#2+.3}) {Fase};}
\newcommand{\efcyl}[5]{% fases, y, nom, valors, color
  \node[left,font=\small\bfseries] at (-.55,#2+.45) {#3};
  \node[left,font=\scriptsize] at (-.05,#2+.9) {+}; \node[left,font=\scriptsize] at (-.05,#2) {$-$};
  \draw[gray!60] (0,#2)--(#1,#2); \draw[gray!60] (0,#2+.9)--(#1,#2+.9);
  \foreach \v [count=\i from 0, remember=\v as \prev] in {#4}{%
    \ifnum\i>0 \draw[#5,line width=1.1pt] ({\i-1},{#2+\prev*.9})--(\i,{#2+\v*.9});\fi}}
\newcommand{\espaifaseu}[3]{%
  \begin{tikzpicture}[x=1.15cm,y=1cm]\efframe{#1}{1.25}\efcyl{#1}{0}{#2}{#3}{vermell}\end{tikzpicture}}
\newcommand{\espaifase}[5]{%
  \begin{tikzpicture}[x=1.15cm,y=1cm]\efframe{#1}{2.95}\efcyl{#1}{1.7}{#2}{#3}{vermell}\efcyl{#1}{0}{#4}{#5}{blau}\end{tikzpicture}}

% Polsador 3/2 normalment tancat amb retorn per molla (cantonada inferior esquerra a (x,y))
\newcommand{\pnbutton}[3]{\pnvalve{#1}{#2}{#3}{32nc}{pol}{spring}}
% Etiqueta d'un component en un esquema: \pntag{x}{y}{text}
\newcommand{\pntag}[3]{\node[pn tag] at (#1,#2) {#3};}
